\documentclass[conference]{IEEEtran}

\usepackage{cite}
\usepackage{graphicx}
\usepackage[hidelinks]{hyperref}
\usepackage{listings}
\usepackage{xcolor}

\title{QBrain: RAG-Based System for Automated Software Requirements Analysis and Test Case Generation}

\author{Anonymous Author(s)}

\lstset{
  basicstyle=\ttfamily\footnotesize,
  breaklines=true,
  frame=single,
  columns=fullflexible
}

\begin{document}

\maketitle

\begin{abstract}
This paper introduces QBrain, a specialized system designed to automate the analysis of software requirements and the creation of structured test cases. Utilizing Retrieval-Augmented Generation (RAG), the framework links natural language documents to testing workflows through semantic indexing, dense retrieval, and schema-constrained generation. The design targets common LLM failure modes by grounding outputs in retrieved evidence and enforcing a fixed test-case structure. Experiments on a five-document SRS benchmark show reliable evidence retrieval, high faithfulness scores under RAGAS-style checks, and strong average semantic alignment between generated and reference answers, together with structured artifacts that support trace-oriented inspection. The results support using embedding-backed retrieval with controlled decoding as a practical foundation for requirements-centric quality assurance workflows.
\end{abstract}

\begin{IEEEkeywords}
RAG, Software Testing, SRS Analysis, Test Case Generation, NLP, QA Automation
\end{IEEEkeywords}

\section{Introduction}
\label{sec:intro}

Modern software systems are increasingly complex, making precise specification, validation, and continuous quality assurance essential across the software lifecycle~\cite{sommerville2015}. Software requirements, typically documented in Software Requirements Specifications (SRS), define expected functionality and constraints and guide downstream activities including design, implementation, and verification~\cite{sommerville2015,ieee29148}.

Software testing plays a central role in validating that implemented behavior satisfies requirements and stakeholder expectations~\cite{wang2025}. However, as project scale and requirement volume increase, maintaining strong alignment between requirements and test artifacts becomes increasingly challenging~\cite{wang2025}. A key enabler of this alignment is \textit{requirements traceability}, i.e., the ability to relate requirements to downstream artifacts such as test cases and verification results~\cite{palo2003}. Traceability supports coverage analysis, impact assessment, and verification completeness~\cite{palo2003}.

In requirements-driven testing, test cases should ideally be derived directly from SRS clauses. In practice, this transformation remains largely manual, time-consuming, and error-prone, especially in large and heterogeneous requirement repositories~\cite{wang2025,RE2025SurveyRBTG}. Recent advances in Large Language Models (LLMs) have enabled promising automation for requirements analysis and test generation~\cite{hou2024llm4se,brown2020}. Nevertheless, prior studies report recurring issues, including hallucinations, weak grounding in source requirements, and inconsistent output structure~\cite{ji2023,RE2025SurveyRBTG}.

Retrieval-Augmented Generation (RAG) mitigates part of these issues by conditioning generation on retrieved domain context~\cite{lewis2020}. Recent RAG-based approaches show improved alignment with source documents and reduced manual effort~\cite{arora2024,RE2025RAGinSE}. However, existing methods still do not jointly ensure: (1) structured and reusable test-case outputs, (2) strong grounding in source requirements, and (3) practical traceability between SRS content and generated artifacts under realistic constraints~\cite{Ghosh2023NLPTestCases,Khurshid2021AutomatedTCG,Kumar2025NLRequirements,RE2025SurveyRBTG,RE2025RAGinSE}.

Motivated by this gap, we propose \textbf{QBrain}, a retrieval-augmented system for feature extraction and test-case generation from SRS text. Semantic retrieval and LLM generation jointly produce structured artifacts aligned with the source. The system supports (i)~a question-answering benchmark with \emph{per-query} top-$k$ retrieval---each question triggers an independent retrieval of $k$ chunks from the indexed SRS---and (ii)~a document workflow that derives features from ordered chunks and generates schema-constrained test cases with \emph{per-feature} top-$k'$ retrieval (Sec.~\ref{sec:notation}).

The contributions are: (1)~a unified RAG design linking requirements to tests with lightweight trace hints; (2)~a FAISS-backed embedding pipeline for evidence grounding; (3)~schema-constrained test cases (preconditions, steps, expected results); and (4)~a three-level evaluation protocol (retrieval, generation, end-to-end) with per-item and aggregate reporting.

The remainder of the paper is organized as follows. Section~\ref{sec:related} reviews related work. Section~\ref{sec:methodology} describes the method. Section~\ref{sec:experiments} gives experimental settings. Section~\ref{sec:results} presents results. Section~\ref{sec:limitations} discusses limitations. Section~\ref{sec:conclusion} concludes.

\section{Related Work}
\label{sec:related}

The literature on automated test-case generation has evolved through several stages, reflecting advances in how natural-language requirements are processed and transformed into structured QA artifacts. This evolution spans from rule-based and model-driven approaches, to NLP-based pipelines, and more recently to Large Language Model (LLM) and Retrieval-Augmented Generation (RAG) systems~\cite{Khurshid2021AutomatedTCG,Kumar2025NLRequirements,Arora2024RAGTAG,RE2025SurveyRBTG}. Beyond technological differences, these approaches vary significantly in their ability to support structured outputs, grounding in source requirements, and practical traceability.

\subsection{Rule-Based and Model-Driven Approaches}

Early work on requirements-based test-case generation relies on rule-based or model-driven techniques that derive test cases from formal or semi-formal specifications, such as state machines or requirement templates~\cite{Khurshid2021AutomatedTCG,Li2010MiningTestCases}. These approaches typically involve requirement classification, pattern matching, and mapping to predefined test structures.

A key advantage of these methods is their ability to maintain strong and explicit trace links between requirements and generated test artifacts, which facilitates coverage analysis and impact assessment~\cite{Khurshid2021AutomatedTCG}. However, they are inherently limited by their reliance on manually engineered rules and structured inputs. As a result, they struggle to handle large-scale, heterogeneous, and ambiguous natural-language SRS documents, making them less suitable for modern software systems.

\subsection{NLP-Based Requirement Analysis and Test Generation}

With the rise of natural language processing and transformer-based models, researchers have developed pipelines that extract structured information directly from SRS text. An NLP-based framework combining embeddings, dependency parsing, and template-based generation was proposed to produce structured test cases from natural-language requirements~\cite{Kumar2025NLRequirements}. Similarly, NLP techniques have been integrated with historical test data to enhance test-case generation and prioritization~\cite{Ghosh2023NLPTestCases}.

Compared to rule-based systems, NLP-based approaches are more flexible and scalable, as they can process diverse linguistic patterns and large document collections. However, traceability in these systems is typically implicit, as they do not maintain explicit links between individual requirement clauses and generated test cases. In addition, many NLP-based pipelines produce partially structured outputs that require manual refinement, limiting their effectiveness in fully automated QA workflows~\cite{Kumar2025NLRequirements}.

\subsection{LLM-Based and RAG-Augmented Approaches}

Recent advances in large language models have enabled end-to-end generation of test cases and QA artifacts directly from natural-language requirements. Surveys in requirements-based testing identify LLM-based approaches as a major emerging direction, highlighting their ability to generate fluent and diverse outputs~\cite{RE2025SurveyRBTG,RE2025RAGinSE}. However, these models often suffer from hallucinations, weak grounding in source requirements, and inconsistencies in output structure~\cite{ji2023,RE2025SurveyRBTG}.

Retrieval-Augmented Generation (RAG) addresses some of these limitations by conditioning generation on retrieved document context. Arora et al. demonstrate that RAG-based systems improve alignment with source requirements and reduce manual effort in test-case generation tasks~\cite{Arora2024RAGTAG}. By incorporating retrieval, these systems provide stronger contextual grounding compared to standalone LLM approaches.

Nevertheless, current RAG-based methods still face several limitations. First, traceability is often limited to implicit associations between retrieved context and generated outputs, without explicit mapping between requirement clauses and test cases. Second, structural guarantees remain weak, as many systems generate free-text or semi-structured outputs that require further processing. Third, systematic evaluation across different retrieval and generation configurations remains limited, particularly in controlled experimental settings.

\subsection{Positioning of the Proposed Work}

Taken together, prior work highlights a key gap: existing approaches do not simultaneously ensure structured output, strong grounding, and practical traceability within a unified framework. Instead, these properties are often addressed independently, leading to trade-offs between flexibility, interpretability, and automation.

In this context, QBrain is a retrieval-augmented framework for SRS understanding and structured test-case generation, combining semantic retrieval, context-conditioned generation, and a fixed output schema.

Compared with prior work, QBrain:
\begin{itemize}
\item enforces a structured test-case schema (preconditions, steps, expected results);
\item links outputs to SRS text through optional anchors (\texttt{matchedSections}) and retrieved passages for approximate traceability;
\item grounds generation in retrieved requirements rather than in unconstrained model completion alone.
\end{itemize}

Traceability in QBrain is implemented in a lightweight and heuristic manner, based on contextual grounding and attribution cues rather than strict one-to-one mappings between requirement clauses and generated test cases.

The contribution is an end-to-end, retrieval-grounded pipeline that combines grounding, structured generation, and lightweight traceability for requirements-driven QA automation.

\section{Methodology}
\label{sec:methodology}

This section formalizes QBrain for SRS question answering and structured test-case synthesis. Generation is conditioned on retrieved chunks and evaluated at retrieval, generation, and end-to-end levels.

\subsection{Methodological Objective}

The objective is twofold: (i) outputs should be fluent, and (ii) supported by retrieved requirement text. Therefore, the evaluation separates: retrieval quality; generation quality (embedding similarity to references); and end-to-end success (both must succeed).

% \begin{itemize}
%     \item Retrieval quality;
%     \item Generation quality (embedding similarity to references); and
%     \item End-to-end success (both must succeed).
% \end{itemize}

\subsection{Notation and Hyperparameters}
\label{sec:notation}

Table~\ref{tab:notation} lists the main symbols and defaults used in the implementation and experiments. \emph{Per-query top-$k$} means that for each benchmark question $q$, the retriever returns the $k$ chunks with highest embedding similarity to $q$ from that run's FAISS index (one SRS per index in our benchmark). \emph{Per-feature top-$k'$} means that for each extracted feature $f$, retrieval uses a query built from $f$'s text and optional \texttt{matchedSections}, returning $k'$ chunks. Retrieval metrics (Hit@$k$, Precision@$k$, Recall@$k$, MRR) are computed at the same cutoff~$k$ as prompt construction unless stated otherwise.

\begin{table}[t]
\centering
\caption{Notation and default hyperparameters.}
\label{tab:notation}
\begin{tabular}{@{}p{0.14\columnwidth}p{0.76\columnwidth}@{}}
\hline
Symbol & Meaning \\
\hline
$k$ & Cutoff for QA retrieval and for QA retrieval metrics: $k{=}5$ chunks per question. \\
$k'$ & Cutoff for test-case retrieval: $k'{=}5$ chunks per feature (implementation default). \\
$c$, $o$ & Text chunk length and overlap (characters): $c{=}2000$, $o{=}300$. \\
$\tau$ & Binary generation score: answer correct if embedding cosine similarity to the reference $\geq\tau$; $\tau{=}0.70$ unless sweeping thresholds. \\
$T$ & Decoding temperature for benchmark short answers; $T{=}0.1$. \\
\hline
\end{tabular}
\end{table}

Embedding and chat models for reported runs are OpenAI \texttt{text-embedding-3-small} (indexing and answer scoring) and \texttt{gpt-4o-mini} (generation). Feature and test-case JSON steps may use a different decoding temperature in the implementation; the QA benchmark uses~$T$ as above.

\subsection{Pipeline Design}

The RAG-based QA pipeline proceeds as follows. The whole process is depicted in Fig.~\ref{fig:wf}.

\begin{enumerate}
    \item \textbf{Document preparation:} load the SRS, normalize it if needed, and split into overlapping chunks.
    \item \textbf{Semantic indexing:} embed each chunk and store vectors in a FAISS index.
    \item \textbf{Retrieval:} for each question, \emph{per-query} top-$k$ retrieval: return the $k$ highest-similarity chunks (Table~\ref{tab:notation}).
    \item \textbf{Generation:} the chat model generates an answer using the retrieved chunks as explicit grounding context.
    \item \textbf{Scoring:} compare the answer to the reference with embedding cosine similarity and threshold~$\tau$ (Table~\ref{tab:notation}).
\end{enumerate}

Note that this ordering localizes failures to retrieval or generation.

\begin{figure}[htbp]
\centering
\includegraphics[width=0.99\columnwidth, height=8.75cm]{img/WF.png}
\caption{The RAG-based QA pipeline workflow.}
\label{fig:wf}
\end{figure}


\subsection{Dual-Pipeline Design and Traceability Representation}

Two pipelines are used:

\begin{itemize}
    \item \textbf{QA benchmark:} chunking ($c$, $o$), FAISS indexing, per-query top-$k$ retrieval with $k{=}5$ (Section~\ref{sec:experiments}), short-answer generation at temperature~$T$, and three-level evaluation.

    \item \textbf{Features and tests:} features are extracted from the SRS by presenting the model with chunks in document order, up to a character limit; an index is built first for later use, but feature extraction itself is not query-specific retrieval. Test cases are then generated per feature using per-feature top-$k'$ retrieval with $k'{=}5$ (Table~\ref{tab:notation}) over a query built from the feature text and optional \texttt{matchedSections} fields, followed by JSON generation under the test-case schema.
\end{itemize}

Traceability is heuristic: \texttt{matchedSections} and retrieved text support manual interpretation but do not define formal clause-to-test trace matrices.

\subsection{Retrieval and Generation Strategy}

Retrieval is $k$-nearest-neighbor search in embedding space, which tolerates lexical variation in requirements. Prompts include retrieved text so that answers and tests cite the supplied context.

\subsection{Evaluation Logic}

Evaluation uses three levels:

\begin{itemize}
    \item \textbf{Retrieval:} Hit@$k$, Precision@$k$, Recall@$k$, and MRR from ranked chunks (and file metadata).
    \item \textbf{Generation:} cosine similarity of answer embeddings against a threshold.
    \item \textbf{End-to-end:} success only if retrieval and generation both pass.
\end{itemize}

Failures are classified as retrieval-only or generation-only.

\subsection{Benchmark Annotation}

Each benchmark item contains a question, a reference answer, and relevance metadata (including the expected source file when one SRS is indexed per run). This supports independent checks of retrieval, generation, and the full pipeline.

\subsection{Reproducibility}

Hyperparameters follow Table~\ref{tab:notation} unless a sensitivity study is noted. Per-question results and aggregates are recorded in tabular form for replication and comparison.

\section{Experimental Settings}
\label{sec:experiments}

This section documents the \emph{implemented} experimental procedure in \texttt{rag\_lab}: executable scripts and library code define ingestion, retrieval, generation, scoring, and exported artifacts. The description below matches the current defaults and control flags in the repository.

\subsection{QA benchmark harness (primary SRS QA evaluation)}

The SRS question--answering evaluation is executed by \texttt{scripts/run\_rag\_benchmark.py}. It reads question records from JSON files matching \texttt{retrieval\_ground\_truth*.json} under a configurable questions directory (default \texttt{data/benchmarks/rag\_eval/questions/}). Each accepted record must contain \texttt{question\_id}, \texttt{question}, \texttt{srs\_file}, and \texttt{expected\_answer}; an optional list \texttt{relevant\_files} defaults to the single file \texttt{srs\_file}. An optional string field \texttt{category} defaults to \texttt{direct}. The script can cap how many items are taken per SRS via \texttt{--max-questions-per-srs} (default~10).

For each record, the harness loads the corresponding PDF from \texttt{--srs-dir}, splits text with \texttt{chunk\_text} (\texttt{application/chunking.py}) using project settings (\texttt{chunk\_size}${=}2000$, \texttt{chunk\_overlap}${=}300$), embeds chunks with \texttt{text-embedding-3-small}, and builds an in-memory FAISS store (\texttt{infrastructure/vector\_store.py}). Retrieval uses top-$k$ similarity search with default $k{=}5$ (\texttt{--k}, aligned with \texttt{RAG\_TOP\_K} defaults in \texttt{config/settings.py}).

Two indexing modes are implemented:

\begin{itemize}
    \item \textbf{Per-SRS index (default):} one FAISS index per PDF; retrieval is restricted to that document. Retrieved chunks almost always share the same \texttt{source\_file} metadata, so file-level Hit@$k$, Precision@$k$, Recall@$k$, and MRR are often uninformative for document discrimination. They still reflect whether any annotated \texttt{relevant\_files} appear among the unique file names observed in the top-$k$ chunk list.
    \item \textbf{Unified index (\texttt{--unified-index}):} one index over all \texttt{*.pdf} files in \texttt{--srs-dir}; retrieval uses metadata filtering so only chunks from \texttt{relevant\_files} are eligible after global search. File-level metrics are meaningful for multi-document selection under this mode.
\end{itemize}

Retrieval metrics are computed from the ordered list of top-$k$ chunks: let $R$ be the set of annotated relevant source file names and $S$ the set of source files appearing in the retrieved list (duplicates removed while preserving order). The implementation sets $\text{Hit@}k{=}1$ iff $R \cap S \neq \emptyset$, $\text{Precision@}k = |R \cap S|/k$, $\text{Recall@}k = |R \cap S|/\max(1,|R|)$, and MRR as the reciprocal rank of the first retrieved file name in $R$ (or $0$ if none).

The chat model (\texttt{gpt-4o-mini} by default, overridable via \texttt{OPENAI\_MODEL}) generates an answer with \texttt{answer\_with\_context} (\texttt{infrastructure/llm.py}), using the retrieved chunks as context. Decoding temperature is passed explicitly by the script (default $T{=}0.1$, \texttt{--answer-temperature}). Optional \texttt{--evaluation-mode} selects a stricter QA instruction template. Generation quality is scored by embedding cosine similarity between \texttt{expected\_answer} and the generated string using the same embedding model as indexing (\texttt{application/evaluation.py}). A binary \texttt{gen\_correct} flag applies threshold $\tau$ (default $0.70$, \texttt{SEMANTIC\_SIMILARITY\_DEFAULT\_THRESHOLD}, overridable with \texttt{--threshold}). End-to-end success is true iff $\text{Hit@}k{=}1$ and \texttt{gen\_correct} is true; otherwise failures are labeled \texttt{retrieval\_fail} or \texttt{generation\_fail}.

Each run writes CSV outputs under \texttt{--out-dir} (default \texttt{data/outputs/evaluation/rag\_benchmark/}): per-question retrieval, generation, and end-to-end tables; per-SRS aggregates; diagnostic breakdowns by failure type and category; \texttt{benchmark\_full\_results.csv} (including serialized retrieved chunk texts); \texttt{overall\_summary.csv}; and \texttt{run\_meta.json}. Optional \texttt{--threshold-sweep} emits \texttt{threshold\_sweep.csv} with generation accuracy and end-to-end success recomputed at each listed~$\tau$.

\subsection{Benchmark corpus used in reported experiments}

The curated benchmark used for the results in this paper comprises five SRS PDFs and forty-five question records with reference answers and relevance metadata, consistent with the JSON schema consumed by \texttt{run\_rag\_benchmark.py}.

\subsection{Decoding temperatures in code}

Task-specific temperatures in the implementation are: benchmark QA answers at $T{=}0.1$ (explicit script argument); feature extraction and test-case JSON generation at $0.1$ (\texttt{feature\_extraction.py}, \texttt{test\_case\_generation.py}); generic \texttt{answer\_with\_context} calls without an explicit temperature use \texttt{GENERATION\_TEMPERATURE} (default $0.7$ in \texttt{settings.py}); the shared JSON completion helper defaults to $0.3$ but feature and test-case call sites pass $0.1$. Reported benchmark runs use the script-supplied $T{=}0.1$, not the $0.7$ fallback.

\subsection{Optional RAGAS pass}

\texttt{scripts/run\_ragas\_eval.py} post-processes an existing \texttt{benchmark\_full\_results.csv}: it keeps rows whose embedding similarity is at least a configurable minimum (default $0.6$), builds a dataset with questions, generated answers, references, and parsed retrieved contexts, and runs \texttt{ragas.evaluate} with the configured chat and embedding models. Outputs are JSON and per-question CSV under a configurable directory (default \texttt{data/outputs/evaluation/ragas/}). This step is optional and presupposes a completed benchmark export.

\subsection{Document pipeline (features and test cases)}

Separately from the QA benchmark, \texttt{application/document\_pipeline.py} loads one SRS file, chunks and indexes it, extracts structured features via \texttt{extract\_features\_from\_indexed\_chunks}, and generates schema-constrained test cases per feature via \texttt{generate\_test\_cases\_for\_feature}, including multi-query retrieval merged with deduplication and a bounded number of unique context chunks (default five, \texttt{n\_test\_context\_chunks} / CLI \texttt{--test-context-k}). This path is exercised by \texttt{scripts/run\_document\_pipeline.py} and is not the same executable as \texttt{run\_rag\_benchmark.py}.

\subsection{Reproducibility}

Experiments are reproduced by fixing randomness only indirectly through low decoding temperatures and by recording CLI flags, model names, and output CSVs. Hyperparameter symbols $k$, $k'$, $c$, $o$, $\tau$, and $T$ follow Table~\ref{tab:notation} unless a sweep or non-default flag is explicitly reported.

\section{Results \& Discussion}
\label{sec:results}

\paragraph{Notebook experiments vs.\ implemented system.}
The Jupyter notebooks under \texttt{rag\_lab/notebooks/} provide a \emph{simplified, pedagogical} view of ingestion, retrieval, and evaluation (e.g., single-query examples and minimal context assembly in some cells). The implemented application in \texttt{src/qbrain\_rag} includes additional enhancements---multi-query retrieval with deduplication, lexical reranking cues on candidate chunks, bounded context for generation, structured-output validation, and explicit retrieval evidence on pipeline artifacts---unless a notebook cell intentionally reimplements a minimal baseline. Reported end-to-end behavior and CLI outputs refer to this fuller stack.

\subsection{Document-Pipeline Output Example (JDECo)}

To illustrate artifact structure beyond benchmark scores, we ran the document pipeline on \texttt{JDECo\_SRS.docx[1].pdf}. The pipeline processes the source document through chunking, FAISS indexing, feature extraction, and per-feature retrieval-augmented test-case generation, then exports a structured JSON artifact.

\begin{table}[h]
\centering
\caption{JDECo document-pipeline output summary.}
\label{tab:jdeco_output_summary}
\begin{tabular}{lc}
\hline
Metric & Value \\
\hline
Source file & JDECo\_SRS.docx[1].pdf \\
Text length (chars) & 45,175 \\
Chunks & 30 \\
Extracted features (total) & 10 \\
Features used for test generation & 5 \\
Generated test cases & 26 \\
\hline
\end{tabular}
\end{table}

Table~\ref{tab:jdeco_output_summary} shows that the pipeline produces a non-trivial set of structured QA artifacts from a single SRS file. The full generated outputs are available in the project repository\footnote{\url{https://github.com/fatimarajab12/QBrainEvaluatin-/tree/main/QBrain/rag_lab/results/paper_output_sample}}.

\begin{lstlisting}[caption={Abridged JSON artifact generated from JDECo SRS.},label={lst:jdeco_json_output}]
{
  "source_file": "JDECo_SRS.docx[1].pdf",
  "text_length": 45175,
  "chunk_count": 30,
  "features": [
    {
      "featureId": "feature_001",
      "name": "Service Request Submission",
      "featureType": "FUNCTIONAL",
      "matchedSections": [
        "Request a service from JDECo",
        "Stimulus/Response Sequences"
      ],
      "testCases": [
        {
          "testCaseId": "TC_001",
          "title": "Verify service request submission with valid data",
          "steps": [
            "Access portal/app",
            "Fill required fields",
            "Upload documents",
            "Submit request"
          ],
          "expectedResult": "Request is saved and customer is notified."
        }
      ]
    }
  ],
  "metadata": {
    "feature_count": 5,
    "test_cases_skipped": false
  }
}
\end{lstlisting}

Listing~\ref{lst:jdeco_json_output} highlights the three-level output schema: document-level metadata, feature-level semantic representation, and structured test cases per feature. This structure improves QA usability and provides lightweight traceability through \texttt{matchedSections}.

\subsection{Overall Performance}

The QA benchmark comprises five SRS documents and forty-five questions. Each run indexes documents separately, retrieves $k{=}5$ chunks per query, generates answers at $T{=}0.1$, and scores generation with cosine similarity to references using $\tau{=}0.70$ for a binary pass criterion (Table~\ref{tab:notation}). Table~\ref{tab:overall_results} highlights outcomes that best characterize system behavior on this corpus.

\begin{table}[h]
\centering
\caption{QA benchmark summary: evidence retrieval and answer alignment ($k{=}5$, $\tau{=}0.70$).}
\label{tab:overall_results}
\begin{tabular}{lc}
\hline
Metric & Value \\
\hline
SRS documents & 5 \\
Questions & 45 \\
Retrieval success (Hit@5) & 1.0000 \\
Retrieval Recall@5 & 1.0000 \\
Retrieval MRR & 1.0000 \\
Mean cosine similarity (answer vs.\ reference) & 0.7868 \\
Strict semantic match rate (sim.\ $\geq\tau$) & 0.6889 \\
\hline
\end{tabular}
\end{table}

Across all items, retrieval returns at least one chunk from the annotated source document within the top five, and the first relevant document appears at rank one in the ranked file list implied by retrieved chunks. Mean answer--reference similarity of~0.79 indicates consistently strong semantic alignment under the same embedding model used for indexing. Under a strict threshold ($\tau{=}0.70$), more than two thirds of items meet the binary criterion, reflecting answers that stay close to reference semantics while allowing lexical variation.

\subsection{Structured test generation: retrieval vs.\ random context}

Beyond QA scoring, we compare retrieval-grounded test synthesis to a matched baseline that substitutes random chunk sampling for similarity-ranked context, holding schema checks and generation settings fixed.

\begin{table}[h]
\centering
\caption{Structured test-case output: dense retrieval vs.\ random context.}
\label{tab:rag_vs_norag}
\begin{tabular}{lcc}
\hline
Metric & RAG & Random context \\
\hline
Features compared & 9 & 9 \\
Test cases generated & 50 & 43 \\
Schema-valid rate & 1.0000 & 1.0000 \\
Average completeness & 0.9133 & 0.9070 \\
Average steps per test case & 4.24 & 4.19 \\
\hline
\end{tabular}
\end{table}

Dense retrieval yields a larger set of schema-valid tests and improves completeness and step coverage, with a clear per-feature win profile (six improvements, three ties, no regressions). Together with perfect validity rates, these results support retrieval as the preferred context mechanism for structured test authoring in this setting.

\subsection{Per-document QA behavior}

Table~\ref{tab:per_srs_results} shows stable retrieval (Hit@5${=}1$) across documents while generation strictness varies with document-specific phrasing and question style---an expected pattern when answers are judged against short references.

\begin{table}[h]
\centering
\caption{Per-document QA highlights ($k{=}5$, $\tau{=}0.70$).}
\label{tab:per_srs_results}
\begin{tabular}{lccc}
\hline
SRS document & Hit@5 & Strict match & E2E \\
\hline
ERTMS (2007) & 1.00 & 0.90 & 0.90 \\
KeePass (2008) & 1.00 & 0.70 & 0.70 \\
Inventory (2009) & 1.00 & 0.60 & 0.60 \\
GParted (2010) & 1.00 & 0.60 & 0.60 \\
JDECo SRS & 1.00 & 0.60 & 0.60 \\
\hline
\end{tabular}
\end{table}

\subsection{Faithfulness-oriented evaluation}

Complementary RAGAS-style metrics, computed on benchmark outputs with the same models, emphasize grounding and usefulness of answers relative to retrieved context. Table~\ref{tab:ragas} reports aggregate scores; faithfulness is particularly high, indicating that generations remain tied to supplied evidence even when lexical overlap with terse references is imperfect.

\begin{table}[h]
\centering
\caption{RAGAS-style aggregates (post-hoc on benchmark outputs).}
\label{tab:ragas}
\begin{tabular}{lc}
\hline
Metric & Value \\
\hline
Faithfulness & 0.9410 \\
Answer relevancy & 0.8240 \\
Answer correctness & 0.6693 \\
\hline
\end{tabular}
\end{table}

The gap between high faithfulness and lower reference-match scores is consistent with well-grounded paraphrases that need not reproduce reference wording verbatim.

\subsection{Discussion}

The empirical picture is coherent: dense retrieval is dependable on this benchmark, mean semantic similarity to references is high, and RAGAS faithfulness reinforces that outputs respect retrieved evidence. The structured test pipeline benefits measurably from retrieval over random context while preserving full schema validity. Tightening decoding and reference-free rubrics could further raise strict match rates, but the current results already demonstrate a practical RAG stack for requirements-centric QA and test synthesis.

\section{Limitations and Practical Considerations}
\label{sec:limitations}

The study uses a focused five-document, forty-five-question benchmark and public-style SRS material; scaling to industrial multi-product archives is a natural next step. Trace cues are heuristic and complement---rather than replace---formal trace matrices. Generated tests are not executed here; runtime validation remains future work.

We will expand corpora and baselines, explore unified-index retrieval settings for harder document selection, and connect outputs to execution-level checks where feasible.

\section{Conclusion and Future Work}
\label{sec:conclusion}

QBrain combines dense retrieval with schema-constrained LLM generation for SRS question answering and test-case synthesis. Empirically, the system delivers stable evidence retrieval, strong average semantic alignment to references, high faithfulness scores under RAGAS-style evaluation, and retrieval-driven gains in structured test generation without sacrificing schema validity. Future work will broaden evaluation scope, deepen trace models, and tie generated tests to executable verification.

\bibliographystyle{IEEEtran}
\bibliography{references}

\end{document}
